\chapter{Counting and Paradoxes}

\begin{center}\itshape
Three doors; behind one of them, a car. You pick door~1. The host, who knows where the car is,\\
opens door~3 and shows a goat. Should you switch to door~2?
\end{center}

Probability began in 1654 with letters between Pascal and Fermat about a gambling problem: how to split the
stakes of a game interrupted before its end. For almost three centuries it remained a collection of brilliant
calculations without a firm foundation, and its paradoxes, some of which open this book, were taken as signs
that the subject was not quite mathematics. In 1933 Kolmogorov gave it axioms: a probability is a measure of
total mass~$1$. Every paradox of this chapter dissolves once the probability space is written down.

This chapter sets up those axioms and the rules that follow from them, then uses them on five classical
problems: the birthday problem, Monty Hall, Bertrand's chords, the ballot theorem and Buffon's needle. Most
sample spaces here are finite, and probability is counting; Buffon's needle and Bertrand's chords use areas,
which Chapter~5 puts on a rigorous footing.

\section{Probability spaces}

Volume~0 (Chapter~8) defined a probability on a finite set. To speak of infinitely many coin tosses, or of a point
chosen at random in a disc, one needs infinite sample spaces, and one must decide which subsets are events.

\begin{definition}
Let $\Omega$ be a set. A \emph{$\sigma$-algebra} on $\Omega$ is a family $\mathcal F$ of subsets of $\Omega$ that
contains $\Omega$, is stable under complement, and is stable under countable unions. A \emph{probability} on
$(\Omega,\mathcal F)$ is a map $P:\mathcal F\to[0,1]$ with $P(\Omega)=1$ that is \emph{countably additive}: for
pairwise disjoint $A_1,A_2,\ldots\in\mathcal F$,
\[
P\Bigl(\bigcup_{n\ge1}A_n\Bigr)=\sum_{n\ge1}P(A_n).
\]
The triple $(\Omega,\mathcal F,P)$ is a \emph{probability space}; the elements of $\mathcal F$ are \emph{events}.
\end{definition}

A $\sigma$-algebra is also stable under countable intersections, since $\bigcap A_n$ is the complement of
$\bigcup A_n^c$, and it contains $\varnothing=\Omega^c$. When $\Omega$ is countable, one takes
$\mathcal F=\mathcal P(\Omega)$, all subsets. For uncountable $\Omega$ such as $[0,1]$, Chapter~5 shows that no
reasonable uniform probability can be defined on all subsets: this is why the axioms name $\mathcal F$.

\begin{proposition}\label{prop:rules}
On a probability space, for events $A$, $B$, $A_n$:
\begin{enumerate}[label=(\alph*)]
\item $P(\varnothing)=0$, and $P(A\cup B)=P(A)+P(B)$ if $A\cap B=\varnothing$;
\item $P(A^c)=1-P(A)$, and $A\subseteq B$ implies $P(A)\le P(B)$;
\item $P(A\cup B)=P(A)+P(B)-P(A\cap B)$;
\item \emph{union bound}: $P\bigl(\bigcup_nA_n\bigr)\le\sum_nP(A_n)$, for finitely or countably many events;
\item \emph{continuity}: if $A_1\subseteq A_2\subseteq\cdots$ then $P\bigl(\bigcup_nA_n\bigr)=\lim_nP(A_n)$; if
$A_1\supseteq A_2\supseteq\cdots$ then $P\bigl(\bigcap_nA_n\bigr)=\lim_nP(A_n)$.
\end{enumerate}
\end{proposition}
\begin{proof}
(a) Countable additivity with all $A_n=\varnothing$ gives $P(\varnothing)=\sum_nP(\varnothing)$, a finite
number, so $P(\varnothing)=0$. Then take $A_1=A$, $A_2=B$ and $A_n=\varnothing$ for $n\ge3$.

(b) $A$ and $A^c$ are disjoint with union $\Omega$, so $P(A)+P(A^c)=1$. If $A\subseteq B$, then $B$ is the
disjoint union of $A$ and $B\cap A^c$, so $P(B)=P(A)+P(B\cap A^c)\ge P(A)$.

(c) $A\cup B$ is the disjoint union of $A$ and $B\cap A^c$, and $B$ is the disjoint union of $B\cap A$ and
$B\cap A^c$. Subtracting the two instances of (a) gives $P(A\cup B)-P(B)=P(A)-P(A\cap B)$.

(d) Let $B_1=A_1$ and $B_n=A_n\cap(A_1\cup\dots\cup A_{n-1})^c$. The $B_n$ are disjoint events, $B_n\subseteq A_n$,
and $\bigcup B_n=\bigcup A_n$ (a point of the union lies in $B_n$ for the first $n$ with $\omega\in A_n$). So
$P(\bigcup A_n)=\sum P(B_n)\le\sum P(A_n)$ by (b).

(e) For an increasing sequence, the $B_n$ of (d) are $A_n\cap A_{n-1}^c$, and $A_N=B_1\cup\dots\cup B_N$. So
$P(\bigcup A_n)=\sum_{n\ge1}P(B_n)=\lim_N\sum_{n\le N}P(B_n)=\lim_NP(A_N)$. For a decreasing sequence, apply this
to the increasing sequence $A_n^c$ and use (b).
\end{proof}

The simplest probability spaces are discrete. A series of non-negative terms can be summed in any order and grouped
in any way without changing its value, since every such sum is the supremum of the finite partial sums (Volume~4,
Chapter~2, writes out the proof); this is all one needs.

\begin{proposition}\label{prop:discrete}
Let $\Omega$ be countable and $p:\Omega\to[0,1]$ with $\sum_{\omega\in\Omega}p(\omega)=1$. Then
$P(A)=\sum_{\omega\in A}p(\omega)$ defines a probability on $(\Omega,\mathcal P(\Omega))$, and every probability on
$(\Omega,\mathcal P(\Omega))$ is of this form, with $p(\omega)=P(\{\omega\})$.
\end{proposition}
\begin{proof}
$P(A)$ is a sum of non-negative terms bounded by $1$, and $P(\Omega)=1$. If $A=\bigcup A_n$ with disjoint $A_n$,
the family $(p(\omega))_{\omega\in A}$ is the union of the disjoint families $(p(\omega))_{\omega\in A_n}$, and
summing a non-negative family by packets gives $\sum_{\omega\in A}p(\omega)=\sum_n\sum_{\omega\in A_n}p(\omega)$.
Conversely, a countable set $A$ is the disjoint countable union of its singletons, so countable additivity
forces $P(A)=\sum_{\omega\in A}P(\{\omega\})$.
\end{proof}

\begin{remark}[infinitely many coin tosses]
A sequence of fair coin tosses $(\omega_1,\omega_2,\ldots)\in\{0,1\}^{\N^*}$ cannot live on a countable space: every
single infinite sequence has probability at most $2^{-n}$ for every $n$ (its first $n$ tosses are prescribed),
hence probability $0$, and a countable union of null events is null. So the space of infinite sequences is
uncountable and the probability is not a sum over points. That a probability space carrying infinitely many
independent fair coins exists is \emph{admitted (proved in Volume~6)}; it is equivalent to the existence of the
uniform probability on $[0,1]$, through binary expansions.
\end{remark}

\section{Conditioning and independence}

\begin{definition}
If $P(B)>0$, the \emph{conditional probability} of $A$ given $B$ is $P(A\mid B)=P(A\cap B)/P(B)$. Events
$(A_i)_{i\in I}$ are \emph{independent} if $P(A_{i_1}\cap\dots\cap A_{i_k})=P(A_{i_1})\cdots P(A_{i_k})$ for every
finite set $\{i_1,\dots,i_k\}$ of distinct indices.
\end{definition}

\begin{proposition}\label{prop:bayes}
Let $B$ be an event with $P(B)>0$.
\begin{enumerate}[label=(\alph*)]
\item $A\mapsto P(A\mid B)$ is a probability on $(\Omega,\mathcal F)$.
\item \emph{Chain rule}: $P(A_1\cap\dots\cap A_n)=P(A_1)P(A_2\mid A_1)\cdots P(A_n\mid A_1\cap\dots\cap A_{n-1})$ when
$P(A_1\cap\dots\cap A_{n-1})>0$.
\item \emph{Total probability}: if $(B_k)$ is a finite or countable partition of $\Omega$ into events with
$P(B_k)>0$, then $P(A)=\sum_kP(A\mid B_k)P(B_k)$.
\item \emph{Bayes}: in the situation of (c), $P(B_j\mid A)=\dfrac{P(A\mid B_j)P(B_j)}{\sum_kP(A\mid B_k)P(B_k)}$ when $P(A)>0$.
\end{enumerate}
\end{proposition}
\begin{proof}
(a) $P(\Omega\mid B)=1$, and if the $A_n$ are disjoint, so are the $A_n\cap B$, whose union is
$(\bigcup A_n)\cap B$; divide countable additivity by $P(B)$.
(b) The right-hand side telescopes: each factor $P(A_k\mid A_1\cap\dots\cap A_{k-1})$ is a quotient whose
denominator cancels the previous numerator.
(c) $A$ is the disjoint union of the $A\cap B_k$, and $P(A\cap B_k)=P(A\mid B_k)P(B_k)$.
(d) $P(B_j\mid A)=P(A\cap B_j)/P(A)$; write the numerator as $P(A\mid B_j)P(B_j)$ and the denominator by (c).
\end{proof}

\begin{think}
A disease affects $1\%$ of the population. A test detects it in $99\%$ of sick people, and gives a false alarm
for $5\%$ of healthy people. You test positive. Guess the probability that you are sick before reading on.
\end{think}

\begin{example}\label{ex:medical}
Let $S$ be ``sick'' and $T$ ``the test is positive''. By Bayes,
\[
\begin{aligned}P(S\mid T)&=\frac{0.99\times0.01}{0.99\times0.01+0.05\times0.99}\eqbr=\frac{0.0099}{0.0594}=\frac16.\end{aligned}
\]
Five positive results out of six are false alarms. The reason is visible in a population of $10\,000$: $100$ are
sick and $99$ of them test positive, but $9\,900$ are healthy and $495$ of them test positive too. A rare
condition makes even a good test unreliable on its own; Exercise~\ref{ex:twotests} shows what a second test
does.
\end{example}

Independence of a family is stronger than independence of each pair. Toss two fair coins and let $A$ be ``the
first is heads'', $B$ ``the second is heads'', $C$ ``the two coins agree''. Each has probability $\frac12$, each
pair has an intersection of probability $\frac14$ (for instance $A\cap C$ is ``two heads''), but
$P(A\cap B\cap C)=\frac14\ne\frac18$. Exercise~\ref{ex:pairwise} turns this example into a tool of computer
science.

\begin{proposition}\label{prop:indepcomp}
If $(A_i)_{i\in I}$ are independent, so are the events obtained by replacing any of the $A_i$ by their
complements.
\end{proposition}
\begin{proof}
It suffices to replace one event, say $A_1$, and to check the product formula for a finite family
$A_1^c,A_2,\dots,A_k$ (families not containing $A_1^c$ are unchanged). With $D=A_2\cap\dots\cap A_k$,
$P(A_1^c\cap D)=P(D)-P(A_1\cap D)=P(A_2)\cdots P(A_k)\bigl(1-P(A_1)\bigr)=P(A_1^c)P(A_2)\cdots P(A_k)$. Replacing
the events one at a time handles any number of them.
\end{proof}

\section{Birthdays and collisions}

\begin{think}
How many people must be in a room for the probability that two share a birthday to exceed $\frac12$? And how many
for a probability that someone shares \emph{your} birthday to exceed $\frac12$?
\end{think}

Assume that birthdays are independent and uniform among $N=365$ days. More generally, throw $n$ balls
independently and uniformly into $N$ boxes: this is also a hash function that sends $n$ keys to $N$ slots, and a
\emph{collision} is two balls in the same box.

\begin{proposition}\label{prop:birthday}
Let $q_n$ be the probability that there is no collision among $n$ balls in $N$ boxes. Then
\[
q_n=\prod_{k=1}^{n-1}\Bigl(1-\frac kN\Bigr)\le\exp\Bigl(-\frac{n(n-1)}{2N}\Bigr),
\]
and, if $n(n-1)\le2N$,
\[
\Bigl(1-\frac1e\Bigr)\frac{n(n-1)}{2N}\le1-q_n\le\frac{n(n-1)}{2N}.
\]
\end{proposition}
\begin{proof}
The sample space is $\{1,\dots,N\}^n$ with the uniform probability. The outcomes without collision are the
injective maps, $N(N-1)\cdots(N-n+1)$ of them; dividing by $N^n$ gives the product. Since $1-x\le e^{-x}$ for all
real $x$ (the exponential is convex and $1-x$ is its tangent at $0$), $q_n\le\exp(-\sum_{k<n}k/N)$, and
$\sum_{k=1}^{n-1}k=\frac{n(n-1)}2$.

Upper bound on $1-q_n$: a collision means that some pair $\{i,j\}$ of balls shares a box, an event of probability
$\frac1N$ (whatever box $i$ lands in, $j$ lands there with probability $\frac1N$). The union bound over the
$\binom n2$ pairs gives $1-q_n\le\binom n2\frac1N$. Lower bound: $1-q_n\ge1-e^{-x}$ with $x=\frac{n(n-1)}{2N}\in[0,1]$,
and $1-e^{-x}\ge(1-e^{-1})x$ on $[0,1]$ because $x\mapsto1-e^{-x}$ is concave, hence above its chord from $0$ to $1$.
\end{proof}

\begin{figure}[ht]
\centering
\fig{ch01-birthday}
\caption{The probability $1-q_n$ of a shared birthday among $n$ people, and the approximation $1-e^{-n(n-1)/730}$.
It passes $\frac12$ at $n=23$ and $99\%$ at $n=57$.}
\label{fig:birthday}
\end{figure}

With $N=365$, $q_{23}=0.4927$, so a shared birthday is more likely than not among $23$ people
(Figure~\ref{fig:birthday}). The surprise comes from counting the wrong thing: $23$ people form $253$ pairs, and
each pair has a chance of $\frac1{365}$. In contrast, the probability that one of $n$ other people shares your own
birthday is $1-(364/365)^n$, which passes $\frac12$ only at $n=253$ (Exercise~\ref{ex:mybirthday}).

\begin{beauty}
Collisions appear after about $\sqrt N$ balls, not $N$. This square root is everywhere in computer science.
A hash table with $N$ slots starts to see collisions after about $\sqrt N$ insertions. A cryptographic hash with
$b$-bit outputs can be broken by the \emph{birthday attack}, which finds two inputs with the same hash after
about $2^{b/2}$ trials: this is why $128$-bit security needs $256$-bit hashes. Pollard's rho method for factoring finds a factor $p$ of $n$ after about $\sqrt p$ steps for the same reason: a random-looking sequence
modulo $p$ repeats after about $\sqrt p$ terms.
\end{beauty}

\section{Monty Hall}

The question that opens this chapter was published in 1975 by Steve Selvin, and became famous in 1990 when the
columnist Marilyn vos Savant answered ``switch'' and received thousands of letters, many from mathematicians,
explaining that she was wrong. She was right; the subtlety is in what the host does.

\begin{proposition}\label{prop:monty}
The car is behind door $C$, uniform in $\{1,2,3\}$. The player picks door $1$. The host opens a door $H$ that is
neither door~$1$ nor door~$C$; when he has two choices (that is, when $C=1$), he picks one uniformly. Then
$P(C=1\mid H=3)=\frac13$ and $P(C=2\mid H=3)=\frac23$: switching doubles the chance of winning.
\end{proposition}
\begin{proof}
By total probability over the three values of $C$, $P(H=3)=\frac13\cdot\frac12+\frac13\cdot1+\frac13\cdot0=\frac12$: if
$C=1$ the host opens door~$3$ half the time, if $C=2$ he must open door~$3$, and if $C=3$ he cannot. Bayes gives
$P(C=1\mid H=3)=\frac{\frac13\cdot\frac12}{\frac12}=\frac13$ and $P(C=2\mid H=3)=\frac{\frac13\cdot1}{\frac12}=\frac23$.
\end{proof}

The answer depends on the host's rule, not only on what is seen. If the host does not know where the car is and
opens door~$2$ or~$3$ at random, then the event ``he opens door~$3$ and it hides a goat'' has probability
$\frac13\cdot\frac12+\frac13\cdot\frac12=\frac13$ (car behind door~$1$ or door~$2$, and he picks door~$3$), and
$P(C=1\mid\text{door 3, goat})=\frac{1/6}{1/3}=\frac12$. Same doors, same goat, different answer: the
information carried by an observation depends on how it was produced.

\section{Bertrand's paradox}

In his \emph{Calcul des probabilités} (1889), Joseph Bertrand asked: what is the probability that a chord drawn
at random in a circle is longer than the side of the inscribed equilateral triangle? He gave three answers.

\begin{proposition}\label{prop:bertrand}
In the unit circle, whose inscribed equilateral triangle has side $\sqrt3$:
\begin{enumerate}[label=(\alph*)]
\item if the two endpoints of the chord are independent and uniform on the circle, $P(\text{length}>\sqrt3)=\frac13$;
\item if the chord is perpendicular to a radius chosen uniformly and crosses it at a uniform distance $r\in[0,1]$
from the centre, the probability is $\frac12$;
\item if the midpoint of the chord is uniform in the disc, the probability is $\frac14$.
\end{enumerate}
\end{proposition}
\begin{proof}
A chord at distance $r$ from the centre has length $2\sqrt{1-r^2}$, which exceeds $\sqrt3$ iff $r<\frac12$.

(a) By rotation, fix the first endpoint at angle $0$; the second is at a uniform angle $\varphi\in[0,2\pi)$, and
the chord has length $2\sin\frac\varphi2$. This exceeds $\sqrt3$ iff $\sin\frac\varphi2>\frac{\sqrt3}2$, that is
$\frac\varphi2\in(\frac\pi3,\frac{2\pi}3)$, an interval of length $\frac{2\pi}3$ out of $2\pi$ for $\varphi$:
probability $\frac13$.

(b) The distance to the centre is $r$, uniform on $[0,1]$, and $r<\frac12$ has probability $\frac12$.

(c) The midpoint is the foot of the perpendicular from the centre, so the distance of the chord to the centre
is the distance of the midpoint to the centre. It is less than $\frac12$ iff the midpoint lies in the disc of
radius $\frac12$, whose area is a quarter of the unit disc's.
\end{proof}

\begin{figure}[ht]
\centering
\fig{ch01-bertrand}
\caption{Sixty random chords for each of Bertrand's three methods; chords longer than the side of the inscribed
triangle (dashed) are orange. The three pictures look different, and they are: they are three different
probabilities on the set of chords.}
\label{fig:bertrand}
\end{figure}

There is no contradiction: ``a chord at random'' does not specify a probability on the set of chords, and the
three methods specify three different ones (Figure~\ref{fig:bertrand}). The paradox is a warning that every
problem in probability starts by naming its probability space.

\begin{remark}
In 1973 E.~T.~Jaynes argued that method~(b) is the natural one: it is the only one whose answer does not change
when the circle is moved or rescaled while the random lines are thrown from far away, as with straws thrown on a
circle drawn on the floor. Experiments with straws agree with $\frac12$.
\end{remark}

\section{The ballot theorem}

\begin{think}
In an election, candidate A gets $a$ votes and B gets $b<a$ votes. The ballots are counted one by one in a random
order. What is the probability that A is strictly ahead during the whole count? Try $a=2$, $b=1$ by hand.
\end{think}

Encode a count as a path: start at height $0$, go up one step for a vote for A and down one for B. After all
votes the path is at height $a-b$. More generally, a \emph{path} from $(0,x)$ to $(n,y)$ is a sequence of heights
$x=s_0,s_1,\dots,s_n=y$ with $s_{k}-s_{k-1}=\pm1$; it has $u$ up-steps and $d$ down-steps with $u+d=n$ and
$u-d=y-x$, so there are $\binom nu$ of them, with $u=\frac{n+y-x}2$ (and none if this is not an integer in
$[0,n]$).

\begin{lemma}[reflection principle]\label{lem:reflection}
Let $x,y>0$. The paths from $(0,x)$ to $(n,y)$ that touch height $0$ are in bijection with all the paths from
$(0,-x)$ to $(n,y)$.
\end{lemma}
\begin{proof}
For a path $s$ from $x$ to $y$ that touches $0$, let $T$ be the first time with $s_T=0$, and reflect the part
before $T$: $s'_k=-s_k$ for $k\le T$, $s'_k=s_k$ for $k\ge T$. This is a path from $-x$ to $y$. Conversely, a
path from $-x<0$ to $y>0$ must pass through $0$; reflecting it before its first visit to $0$ gives back a path
from $x$ to $y$ touching $0$, and the two maps are inverse to each other because the first visit to $0$ is the
same time $T$ for $s$ and $s'$ (before $T$, $s_k\ne0$ iff $s'_k\ne0$).
\end{proof}

\begin{figure}[ht]
\centering
\fig{ch01-reflection}
\caption{The reflection principle: the part of the path before its first visit $T$ to $0$ is reflected, turning
a path from $a$ to $b$ that touches $0$ into a path from $-a$ to $b$.}
\label{fig:reflection}
\end{figure}

\begin{theorem}[Bertrand's ballot theorem, 1887]\label{thm:ballot}
If A receives $a$ votes and B receives $b<a$ votes, and all $\binom{a+b}a$ orders of the ballots are equally likely,
the probability that A is strictly ahead throughout the count is $\dfrac{a-b}{a+b}$.
\end{theorem}
\begin{proof}
Let $n=a+b$. A good count is a path from $(0,0)$ to $(n,a-b)$ with $s_k>0$ for $1\le k\le n$. Its first step is
up, and the rest is a path from $(1,1)$ to $(n,a-b)$ that never touches $0$. Paths from $(1,1)$ to $(n,a-b)$ have
$n-1$ steps with $a-1$ up-steps: there are $\binom{n-1}{a-1}=\binom{n-1}b$ of them. By the reflection principle (shifted
by one unit of time), those that touch $0$ are as many as the paths from $(1,-1)$ to $(n,a-b)$, which have $a$
up-steps and $b-1$ down-steps: $\binom{n-1}{b-1}$ of them. So the number of good counts is
$\binom{n-1}b-\binom{n-1}{b-1}$, and dividing by $\binom nb$,
\[
\frac{\binom{n-1}b}{\binom nb}-\frac{\binom{n-1}{b-1}}{\binom nb}=\frac{n-b}n-\frac bn=\frac{a-b}{a+b}.
\]
\end{proof}

For $a=2$, $b=1$ the orders are AAB, ABA, BAA and only the first is good: $\frac13=\frac{2-1}{2+1}$.
Exercise~\ref{ex:catalan} deduces the Catalan numbers from the theorem, and Chapter~7 uses the reflection
principle again to study random walks.

\section{Buffon's needle}

In 1777 Georges-Louis Leclerc, comte de Buffon, asked for the probability that a needle thrown on a floor of
parallel boards crosses a line between two boards. It was the first problem of \emph{geometric probability}, where
probabilities are ratios of areas. Here we take the rule ``a uniform point in a rectangle falls in a region with
probability equal to the ratio of areas'' as the model; Chapter~5 constructs the uniform probability and shows
that it agrees with the Riemann integral on regions bounded by graphs of continuous functions.

\begin{think}
A needle of length $\ell$ is thrown on lines spaced $d\ge\ell$ apart. The answer contains $\pi$. Where could
$\pi$ come from in a problem about straight lines?
\end{think}

\begin{theorem}[Buffon]\label{thm:buffon}
Let $\ell\le d$. Describe a needle of length $\ell$ by the distance $x\in[0,\frac d2]$ from its midpoint to the
nearest line and by the acute angle $\theta\in[0,\frac\pi2]$ it makes with the lines, and assume that $(\theta,x)$ is
uniform in the rectangle $[0,\frac\pi2]\times[0,\frac d2]$. Then the needle crosses a line with probability
$\dfrac{2\ell}{\pi d}$.
\end{theorem}
\begin{proof}
The half-needle on the side of the nearest line reaches a distance $\frac\ell2\sin\theta$ from the midpoint,
measured perpendicularly to the lines. So the needle crosses that line iff $x\le\frac\ell2\sin\theta$; it cannot
reach the other neighbouring line, at distance $d-x\ge\frac d2\ge\frac\ell2$, except with $\theta=\frac\pi2$ and
$x=\frac d2=\frac\ell2$, a case of zero area. The favourable region lies under the graph of
$\theta\mapsto\frac\ell2\sin\theta$ (Figure~\ref{fig:buffon}), whose area is $\int_0^{\pi/2}\frac\ell2\sin\theta\dd\theta=\frac\ell2$.
The rectangle has area $\frac\pi2\cdot\frac d2$, and the ratio is $\frac{\ell/2}{\pi d/4}=\frac{2\ell}{\pi d}$.
\end{proof}

\begin{figure}[ht]
\centering
\fig{ch01-buffon}
\caption{Left: forty needles with $\ell=0.8$ on lines one unit apart; the crossing ones are orange. Right: in the
rectangle of $(\theta,x)$, the needle crosses iff the point lies under $x=\frac\ell2\sin\theta$.}
\label{fig:buffon}
\end{figure}

So $\pi$ comes from the uniform angle, and throwing needles is a way to measure~$\pi$: if $c$ of $n$ needles
cross, $\frac{2\ell n}{dc}$ estimates $\pi$. It is a poor method: Chapter~6 shows that the error after $n$ throws is
of order $\frac1{\sqrt n}$, so a million throws give about three correct digits. Exercise~\ref{ex:noodle} gives a
second, astonishing proof of Buffon's formula in which no integral appears.

\begin{remark}[a suspicious experiment]
In 1901 the Italian mathematician Mario Lazzarini reported $3408$ throws of a needle with $\ell/d=\frac56$ and
$1808$ crossings, giving $\pi\approx\frac{2\cdot5\cdot3408}{6\cdot1808}=\frac{355}{113}=3.1415929\ldots$, correct to
six decimals. That is far too good: with $3408$ throws the typical error is in the second decimal. The number of
throws was chosen so that $\frac{355}{113}$, a famous approximation of $\pi$, could come out exactly, and the
experiment was presumably stopped at the right moment. Statistics can detect results that are too good to be true.
\end{remark}

\begin{beauty}
Kolmogorov's axioms say nothing about what probability \emph{means}: frequencies in repeated experiments,
degrees of belief, or symmetries of a physical system. They only say how probabilities combine. This modesty is
their strength. The paradoxes of this chapter were not failures of logic but failures to say which probability
was meant, and once it is named, the rest is mathematics: counting in this chapter, integration from Chapter~5 on.
\end{beauty}

\section{Exercises}

\begin{exercise}[$\star$]\label{ex:mybirthday}
How many other people must be in a room for the probability that one of them shares your birthday to exceed
$\frac12$? Compare with the $23$ of the birthday problem and explain the difference.
\end{exercise}

\begin{exercise}[$\star$]\label{ex:twochildren}
A family has two children, each a boy or a girl with probability $\frac12$, independently. (a) Given that at
least one is a boy, what is the probability that both are boys? (b) Given that at least one is a boy born on a
Tuesday (each day of the week having probability $\frac17$, independently of everything else), what is the
probability that both are boys?
\end{exercise}

\begin{exercise}[$\star$]\label{ex:twotests}
In Example~\ref{ex:medical}, you take a second test, whose errors are independent of the first given your state
of health, and it is positive too. What is now the probability that you are sick?
\end{exercise}

\begin{exercise}[$\star$]\label{ex:penney}
A fair coin is tossed until either HHH or THH appears as three consecutive tosses. Show that THH comes first with
probability $\frac78$.
\end{exercise}

\begin{exercise}[$\star\star$]\label{ex:montyn}
There are $n\ge3$ doors and one car. You pick a door; the host, who knows where the car is, opens $k\le n-2$ other
doors hiding goats, chosen uniformly among the possible ones. You then switch to one of the $n-k-1$ other closed
doors, chosen uniformly. Show that you win with probability $\frac{n-1}{n(n-k-1)}$, and that switching is always
better than staying.
\end{exercise}

\begin{exercise}[$\star\star$]\label{ex:pairwise}
Let $X_1,\dots,X_k$ be independent fair bits. For each non-empty $S\subseteq\{1,\dots,k\}$, let $A_S$ be the event
$\sum_{i\in S}X_i\equiv0\pmod2$. Show that the $2^k-1$ events $A_S$ each have probability $\frac12$ and are pairwise
independent, but that no three of the form $A_S,A_T,A_{S\triangle T}$ are independent. (In computer science, this
turns $k$ random bits into $2^k-1$ pairwise independent ones, which is enough for many algorithms.)
\end{exercise}

\begin{exercise}[$\star\star$]\label{ex:catalan}
A path of length $2n$ from $0$ to $0$ is a \emph{Dyck path} if it never goes below $0$. Using the ballot theorem
with $a=n+1$ and $b=n$, show that there are $C_n=\frac1{n+1}\binom{2n}n$ Dyck paths of length $2n$ (the Catalan
numbers).
\end{exercise}

\begin{exercise}[$\star\star$]\label{ex:secretary}
\emph{The secretary problem.} $n$ candidates with distinct values arrive in a uniformly random order. After each
interview you must accept or reject the candidate, for ever, and you win if you accept the best one. For
$2\le r\le n$, the \emph{threshold strategy} rejects the first $r-1$ candidates and accepts the first later one
who is better than all before. (a) Show that it wins with probability $p_r=\frac{r-1}n\sum_{k=r}^n\frac1{k-1}$.
(b) Show that $\max_rp_r\to\frac1e$ as $n\to\infty$, attained for $r\approx n/e$.
\end{exercise}

\begin{exercise}[$\star\star\star$]\label{ex:noodle}
\emph{Buffon's noodle} (Barbier, 1860). Lines are spaced $d$ apart. For a rigid plane curve $\Gamma$ thrown at
random (with a uniform random position and orientation, as in Theorem~\ref{thm:buffon}), let $N(\Gamma)$ be the
number of its intersection points with the lines.
(a) Using only linearity of expectation, show that for a polygonal line, $\mathbb E[N(\Gamma)]=c\cdot\text{length}(\Gamma)$
for a constant $c$ that does not depend on $\Gamma$. (You may admit that $\mathbb E[N]$ of a segment is
an increasing function of its length that is additive under subdivision, which the model makes clear.)
(b) A circle of diameter $d$ meets the lines in exactly $2$ points with probability $1$. Comparing it with
inscribed and circumscribed regular polygons, show that $c=\frac2{\pi d}$, and deduce Buffon's theorem again.
\end{exercise}

\begin{exercise}[$\star\star$]\label{pb:whichroot}
For $n\ge0$, $\binom{1/2}n=\frac1{n!}\cdot\frac12\bigl(\frac12-1\bigr)\cdots\bigl(\frac12-n+1\bigr)$.
\begin{enumerate}[label=(\alph*)]
\item Show that $\binom{1/2}n=(-1)^{n-1}\frac{C_{n-1}}{2^{2n-1}}$ for $n\ge1$, where $C_m=\frac1{m+1}\binom{2m}m$ is the $m$-th
Catalan number.
\item Show that $\sum_{n\ge0}\binom{1/2}n7^n$ converges in $\Q_7$ to a square root of $8$, and say which one.
\item Show that $\sum_{n\ge0}\binom{1/2}n8^n$ converges in $\Q_2$, and that its sum is $-3$.
\end{enumerate}
\end{exercise}

\begin{exercise}[$\star$]\label{pb:nouniform}
\emph{There is no uniform integer.}
\begin{enumerate}[label=(\alph*)]
\item Show that there is no countably additive probability measure on $\N^*$ that gives probability $\frac1m$ to every residue
class modulo every $m$.
\item Show that density is finitely additive on sets that have a density, and translation invariant.
\item Find two sets $A$ and $B$ that both have density $\frac12$ while $A\cap B$ has no density.
\end{enumerate}
\end{exercise}

\begin{exercise}[$\star\star$]\label{pb:bollobas}
\emph{Bollobás's set-pairs inequality (1965).} Let $A_1,\dots,A_m$ and $B_1,\dots,B_m$ be finite sets with $A_i\cap B_i=\varnothing$ for every
$i$ and $A_i\cap B_j\ne\varnothing$ for $i\ne j$. Let $a_i=|A_i|$, $b_i=|B_i|$.
\begin{enumerate}[label=(\alph*)]
\item Show that $\displaystyle\sum_{i=1}^m\binom{a_i+b_i}{a_i}^{-1}\le1$.
\item A graph is \emph{$\tau$-critical} if deleting any edge lowers the minimum size $\tau$ of a vertex cover. Show that a $\tau$-critical
graph with $\tau=t+1$ has at most $\binom{t+2}2$ edges, and that $K_{t+2}$ attains this.
\end{enumerate}
\end{exercise}

\begin{exercise}[$\star\star\star$]\label{pb:granularity}
\emph{Low degree forces few variables.} Let $f:\{-1,1\}^n\to\{-1,1\}$ have degree $d\ge1$.
\begin{enumerate}[label=(\alph*)]
\item Show that every $\hat f(S)$ is an integer multiple of $2^{1-d}$.
\item Deduce that at most $4^{d-1}$ coefficients $\hat f(S)$ are non-zero.
\item Show that a non-zero multilinear polynomial of degree at most $k$ does not vanish on at least a fraction $2^{-k}$ of the
cube.
\item Call $i$ \emph{relevant} if $\operatorname{Inf}_i(f)>0$. Show that $\operatorname{Inf}_i(f)\ge2^{1-d}$ for every relevant $i$, and deduce that $f$ has at most
$d\,2^{d-1}$ relevant variables.
\end{enumerate}
\end{exercise}

\section{Hints}

\begin{hint}{ex:mybirthday}
Each other person misses your birthday with probability $\frac{364}{365}$, independently.
\end{hint}
\begin{hint}{ex:twochildren}
List the equally likely outcomes. In (b), count the pairs (sex, day) for two children: $14^2$ outcomes.
\end{hint}
\begin{hint}{ex:twotests}
Apply Bayes with the likelihoods $0.99^2$ and $0.05^2$, or apply Example~\ref{ex:medical} again with the prior $\frac16$.
\end{hint}
\begin{hint}{ex:penney}
When can HHH appear before THH? Look at the toss just before the first HHH.
\end{hint}
\begin{hint}{ex:montyn}
Condition on whether your first door hides the car.
\end{hint}
\begin{hint}{ex:pairwise}
For $S\ne T$, pick $i\in S\setminus T$ (or $T\setminus S$) and condition on all the other bits.
\end{hint}
\begin{hint}{ex:catalan}
A good count for $a=n+1$, $b=n$ starts with a vote for A; what does the rest look like?
\end{hint}
\begin{hint}{ex:secretary}
The strategy picks the best candidate if the best is at position $k\ge r$ and the best of the first $k-1$ is
among the first $r-1$. Compare the sum with an integral.
\end{hint}
\begin{hint}{ex:noodle}
(a) Each segment of the polygon contributes its own expected number of crossings. (b) A line meets a convex polygon
in at most two points, and if it crosses the circle it crosses a circumscribed polygon twice.
\end{hint}

\begin{hint}{pb:whichroot}
Vandermonde's identity $\sum_k\binom\alpha k\binom\beta{n-k}=\binom{\alpha+\beta}n$ is a polynomial identity in $\alpha,\beta$.
The two square roots of $8$ in $\Q_7$ differ modulo $7$, and those of $9$ in $\Q_2$ differ modulo $8$.
\end{hint}

\begin{hint}{pb:nouniform}
(a) What would the probability of a single integer be? (c) Let $B$ agree with the even numbers on some long blocks and with the
odd numbers on others.
\end{hint}

\begin{hint}{pb:bollobas}
(a) In a random ordering of the ground set, what is the chance that all of $A_i$ comes before all of $B_i$? Can it happen for two
indices? (b) Apply (a) to the edges and to the covers that miss them.
\end{hint}

\begin{hint}{pb:granularity}
(a) Write $g=\frac{1-f}2$ as a polynomial in the variables $y_i=\frac{1-x_i}2\in\{0,1\}$: its coefficients are integers. (c) Induction on $n$:
write $p=x_nq+r$. (d) Apply (c) to the derivative $D_if$, and use $I(f)=\sum|S|\hat f(S)^2\le d$.
\end{hint}

\section{Solutions}

\begin{solution}{ex:mybirthday}
\idea ``Someone shares my birthday'' fixes the day, so there is no square root: we need about $365\log2$ people,
not $\sqrt{365}$.
\writeup The probability that none of $n$ people has your birthday is $(364/365)^n$, which is less than $\frac12$
iff $n>\frac{\log2}{\log(365/364)}=252.65\ldots$; so $n=253$. In the birthday problem every pair of people can
collide, and $23$ people form $\binom{23}2=253$ pairs: the same number, which is no coincidence, since both
situations need about $253$ ``chances of $\frac1{365}$''.
\end{solution}

\begin{solution}{ex:twochildren}
\idea Write the sample space and count; conditioning on ``at least one'' keeps several outcomes, and extra
details change how many.
\writeup (a) The four outcomes BB, BG, GB, GG are equally likely, and ``at least one boy'' keeps three of them,
of which one is BB: probability $\frac13$.
(b) Each child is one of $14$ equally likely (sex, day) pairs, so there are $196$ outcomes. Those where no child is
a Tuesday boy number $13^2=169$, so ``at least one Tuesday boy'' has $27$ outcomes. Among them, both are boys in
$7^2-6^2=13$ outcomes (two boys, not both born on other days). The probability is $\frac{13}{27}$, close to
$\frac12$: the more specific the information about one child, the closer it is to ``the first child is a boy''.
\end{solution}

\begin{solution}{ex:twotests}
\idea Bayes again: the posterior after the first test is the prior for the second.
\writeup With independent errors, $P(T_1\cap T_2\mid S)=0.99^2$ and $P(T_1\cap T_2\mid S^c)=0.05^2$, so
\[
\begin{aligned}P(S\mid T_1\cap T_2)&=\frac{0.9801\times0.01}{0.9801\times0.01+0.0025\times0.99}\\&=\frac{0.009801}{0.012276}\approx0.798.\end{aligned}
\]
Equivalently, the odds of being sick are multiplied by $\frac{0.99}{0.05}=19.8$ at each positive test: from $\frac1{99}$
to $\frac15$ to $3.96$, a probability of $\frac{3.96}{4.96}\approx0.798$.
\end{solution}

\begin{solution}{ex:penney}
\idea Ask what precedes the first HHH.
\writeup If HHH has not appeared in the first three tosses, look at its first appearance, at tosses
$k,k+1,k+2$ with $k\ge2$. Toss $k-1$ is not H, otherwise HHH would have appeared at $k-1$; so it is T, and THH
appeared at tosses $k-1,k,k+1$, before HHH. So HHH wins only if the first three tosses are HHH, which has
probability $\frac18$. THH does appear eventually with probability $1$: the tosses $3m+1,3m+2,3m+3$ for
$m=0,1,\ldots$ are independent triples, each equal to THH with probability $\frac18$, so the probability that none
of the first $M$ of them is THH is $(\frac78)^M\to0$. Hence THH comes first with probability $\frac78$.
\end{solution}

\begin{solution}{ex:montyn}
\idea If your first door hides the car, switching loses; otherwise the car is behind one of the remaining closed
doors.
\writeup With probability $\frac1n$ your door hides the car and switching loses. With probability $\frac{n-1}n$ it
does not; the host opens $k$ goat doors, so the car is behind one of the $n-k-1$ other closed doors, and you pick
it with probability $\frac1{n-k-1}$. So switching wins with probability $\frac{n-1}{n(n-k-1)}$, which is larger
than the $\frac1n$ of staying since $n-1>n-k-1$ for $k\ge1$.
\end{solution}

\begin{solution}{ex:pairwise}
\idea A parity involving a bit that nothing else involves is a fair coin, whatever the other bits are.
\writeup For $S\ne\varnothing$, pick $i\in S$; given all the other bits, $\sum_{j\in S}X_j$ is a fixed number plus
$X_i$, which is even with probability $\frac12$; by total probability, $P(A_S)=\frac12$. For $S\ne T$, say
$i\in S\setminus T$ (otherwise swap them). Given all bits other than $X_i$, the parity on $T$ is determined, and
the parity on $S$ is even with probability $\frac12$; so $P(A_S\cap A_T\mid\text{others})=\frac12\mathbf
1_{A_T}$, and averaging, $P(A_S\cap A_T)=\frac12P(A_T)=\frac14$. Finally, the sum over $S\triangle T$ is congruent
to the sum over $S$ plus the sum over $T$ modulo $2$ (elements of $S\cap T$ are counted twice), so $A_S\cap A_T\subseteq
A_{S\triangle T}$ and $P(A_S\cap A_T\cap A_{S\triangle T})=\frac14\ne\frac18$.
\end{solution}

\begin{solution}{ex:catalan}
\idea A strict lead for A when $a=b+1$ is a Dyck path lifted by one step.
\writeup With $a=n+1$ and $b=n$, a good count starts with an A, after which the path is at $1$ and must stay at
least at $1$ for the remaining $2n$ steps, ending at $1$. Shifting down by $1$, these remainders are exactly the
Dyck paths of length $2n$. By the ballot theorem their number is
\[
\begin{aligned}\frac{a-b}{a+b}\binom{a+b}a&=\frac1{2n+1}\binom{2n+1}{n}\eqbr=\frac1{2n+1}\cdot\frac{(2n+1)!}{n!\,(n+1)!}\\&=\frac{(2n)!}{n!\,(n+1)!}=\frac1{n+1}\binom{2n}n.\end{aligned}
\]
\end{solution}

\begin{solution}{ex:secretary}
\idea Decompose according to the position $k$ of the best candidate; the strategy then succeeds iff it does not
stop before $k$.
\writeup (a) Let the best candidate be at position $k$; this has probability $\frac1n$ for each $k$. If $k<r$ the
strategy rejects it. If $k\ge r$, the strategy accepts it iff it has not stopped at a position $j\in[r,k-1]$, that
is, iff no candidate between $r$ and $k-1$ beats all the earlier ones; this happens iff the best of the first $k-1$
candidates is among the first $r-1$. By symmetry the best of the first $k-1$ is at a uniform position among
them (and this is independent of where the overall best is), so this has probability $\frac{r-1}{k-1}$. Summing,
$p_r=\sum_{k=r}^n\frac1n\cdot\frac{r-1}{k-1}$.
(b) Write $r-1=xn$. Then $\sum_{k=r}^n\frac1{k-1}=\sum_{j=r-1}^{n-1}\frac1j$ lies between $\log\frac n{r-1}$ and
$\log\frac{n-1}{r-2}$ for $r\ge3$ (compare with $\int\frac{\mathrm dt}t$, Volume~0, Chapter~5), and the two bounds differ by at most $\log\frac{r-1}{r-2}\le\frac1{r-2}$; so $p_r=x\log\frac1x+o(1)$
uniformly for $x\in[\frac1{\sqrt n},1]$, while for smaller $x$ both $p_r\le\frac{r-1}n(1+\log n)$ and
$x\log\frac1x$ tend to $0$. The function $x\log\frac1x$ has its maximum $\frac1e$ at $x=\frac1e$, so
$\max_rp_r\to\frac1e$, with $r\approx\frac ne$. (One can show that no strategy does better than the best threshold
strategy, Lindley 1961.)
\end{solution}

\begin{solution}{ex:noodle}
\idea A curve made of pieces crosses the lines as often as its pieces do, so the expected count is additive in
length. The constant is found on one curve where the count is known for sure: a circle of diameter $d$.
\writeup (a) Let $f(\ell)=\mathbb E[N]$ for a segment of length $\ell$. Cutting a segment into two collinear pieces
of lengths $\ell_1,\ell_2$, every intersection point is on one piece (an endpoint shared by both lies on a line with
probability $0$), so $f(\ell_1+\ell_2)=f(\ell_1)+f(\ell_2)$. Then $f(m\ell)=mf(\ell)$ for integers $m\ge1$, and $kf(\frac mk\ell)=f(m\ell)$, so $f(q\ell)=qf(\ell)$ for rational $q>0$; since $f$ is increasing, $f(\ell)=c\ell$ with $c=f(1)$: for irrational $\ell$,
squeeze it between rationals. A polygonal line is thrown as a rigid body, so each of its segments is a thrown
segment with a uniform position and orientation; by linearity of expectation, $\mathbb E[N]$ is the sum over
the segments, $c$ times the total length.
(b) The lines cut the plane into strips of width $d$. A circle of diameter $d$ whose centre is not on a line
meets exactly one line (the one at distance less than $\frac d2$ from the centre, unless the centre is exactly
halfway, a null event) in exactly two points. Let $P_m$ and $Q_m$ be the regular $m$-gons inscribed in and
circumscribed about the circle, thrown together with it. A line meets the convex polygon $P_m$ in at most $2$
points, and only if it meets the circle; so $N(P_m)\le N(\text{circle})=2$ almost surely. A line crossing the
circle crosses $Q_m$ in $2$ points, so $N(Q_m)\ge2$. Taking expectations, $c\cdot\text{per}(P_m)\le2\le c\cdot\text{per}(Q_m)$.
Both perimeters tend to $\pi d$ (they are $md\sin\frac\pi m$ and $md\tan\frac\pi m$), so $c=\frac2{\pi d}$. For a needle of
length $\ell\le d$, $N\in\{0,1\}$, so $P(\text{crossing})=\mathbb E[N]=\frac{2\ell}{\pi d}$.
\end{solution}

\begin{solution}{pb:whichroot}
\idea The binomial series $\sum\binom{1/2}nx^n$ is ``$\sqrt{1+x}$'' because its square is $1+x$ as a formal power series. In
$\Q_7$, the powers $7^n$ are small, so the series converges; in $\R$ it diverges, since $|x|>1$. Which root it picks is decided
by the first term and the size of the others, and the surprise is that in $\Q_2$ the series for $\sqrt9$ lands on $-3$.

\writeup (a) For $n\ge1$, the numerator of $\binom{1/2}n$ is
\[
\frac12\cdot\Bigl(-\frac12\Bigr)\Bigl(-\frac32\Bigr)\cdots\Bigl(-\frac{2n-3}2\Bigr)=(-1)^{n-1}\frac{1\cdot3\cdots(2n-3)}{2^n},
\]
and $1\cdot3\cdots(2n-3)=\frac{(2n-2)!}{2^{n-1}(n-1)!}$. So
\[
\binom{1/2}n=(-1)^{n-1}\frac{(2n-2)!}{2^{2n-1}\,n!\,(n-1)!}=(-1)^{n-1}\frac{C_{n-1}}{2^{2n-1}},
\]
since $C_{n-1}=\frac{(2n-2)!}{n!(n-1)!}$, an integer (Volume~3, Chapter~1).

\emph{The square of the series.} For non-negative integers $\alpha,\beta$, $\sum_{k=0}^n\binom\alpha k\binom\beta{n-k}=
\binom{\alpha+\beta}n$: choose $n$ people among $\alpha$ women and $\beta$ men. Both sides are polynomials in $(\alpha,\beta)$
that agree on $\N^2$, so they are equal for all $\alpha,\beta$ (a polynomial in two variables vanishing on $\N^2$ is zero:
apply Volume~0, Chapter~3 in each variable). With $\alpha=\beta=\frac12$: $\sum_k\binom{1/2}k\binom{1/2}{n-k}=\binom1n$, which is
$1$ for $n=0,1$ and $0$ for $n\ge2$.

\emph{Products in $\Q_p$.} If $a_n\to0$ and $b_n\to0$ in $\Q_p$, then $\sum a_n$ and $\sum b_n$ converge (Volume~1,
Chapter~12), and so does $\sum c_n$ with $c_n=\sum_{k\le n}a_kb_{n-k}$, to the product of the sums. Indeed, let $\beta=\sup|b_n|$
and $\alpha=\sup|a_n|$. Then
$\bigl(\sum_{k\le N}a_k\bigr)\bigl(\sum_{l\le N}b_l\bigr)-\sum_{n\le N}c_n=\sum a_kb_l$ over $k,l\le N$ with $k+l>N$; in each
term $k>\frac N2$ or $l>\frac N2$, so by the ultrametric inequality the difference has absolute value at most
$\max\bigl(\beta\sup_{k>N/2}|a_k|,\alpha\sup_{l>N/2}|b_l|\bigr)\to0$.

(b) The term $\binom{1/2}n7^n$ has $7$-adic absolute value $7^{-n}$ at most, since $C_{n-1}$ is an integer and $2^{2n-1}$ is
a unit in $\Z_7$. So the series converges to some $s\in\Z_7$, and by the two facts above $s^2=1+7=8$. Every term with $n\ge1$
lies in $7\Z_7$, so $s\equiv1\pmod7$. The square roots of $8$ in $\Q_7$ are $\pm s$, and $-s\equiv-1\not\equiv1$: the series
converges to \emph{the} square root of $8$ that is $\equiv1\pmod7$.

(c) Now $\binom{1/2}n8^n=(-1)^{n-1}C_{n-1}2^{3n-(2n-1)}=(-1)^{n-1}C_{n-1}2^{n+1}$, whose $2$-adic absolute value is at most
$2^{-n-1}$. The series converges in $\Z_2$ to some $s$ with $s^2=9$, so $s=\pm3$. Modulo $8$: the terms with $n\ge2$ are
divisible by $2^{n+1}\ge8$, and the first two are $1$ and $\binom{1/2}18=4$. So $s\equiv5\pmod8$, while $3\not\equiv5$ and
$-3\equiv5\pmod8$. Hence $s=-3$.
\tools Volume~1, Chapter~12 (convergence in $\Q_p$); Volume~0, Chapters~1 and~3 (binomial identities, polynomials); Volume~3,
Chapter~1 (Catalan numbers).
\leads The same series can converge in $\R$ and in several $\Q_p$ to different square roots, a first hint that a rational
number ``lives'' in all completions of $\Q$ at once. This point of view leads to the adeles (Volume~7) and to the local--global
principle (Volume~12).
\end{solution}

\begin{solution}{pb:nouniform}
\idea A uniform probability would give every single integer the same mass, which must be $0$, and countable additivity then
kills everything. Densities avoid the contradiction by giving up countable additivity.

\writeup (a) Suppose $P$ is such a measure. A singleton $\{n\}$ lies in the class of $n$ modulo $m$, so $P(\{n\})\le\frac1m$ for
every $m$, hence $P(\{n\})=0$, and $P(\N^*)=\sum_nP(\{n\})=0$, a contradiction.

(b) If $A,B$ are disjoint with densities, $\#((A\cup B)\cap[1,N])=\#(A\cap[1,N])+\#(B\cap[1,N])$, so $A\cup B$ has density
$d(A)+d(B)$. If $A$ has a density, $\#((A+h)\cap[1,N])$ differs from $\#(A\cap[1,N])$ by at most $|h|$, which is $o(N)$.

(c) Let $E$ be the even numbers and $I_j=[4^j,2\cdot4^j[$. Let $B$ be the set of even numbers in $\bigcup_jI_j$ together with the
odd numbers outside it. In any interval, $B$ contains one of every two consecutive integers, so $B$ has density $\frac12$, and
so does $E$. But $E\cap B$ is the set of even numbers of $\bigcup I_j$. The blocks $I_j$ make up a proportion of $[1,N]$ that is
about $\frac23$ at $N=2\cdot4^j$ and about $\frac13$ at $N=4^{j+1}$, so the proportion of $E\cap B$ oscillates between about
$\frac13$ and $\frac16$, and has no limit.

\tools Volume~3, Chapter~1 (additivity of probabilities).
\leads Densities are finitely additive ``probabilities'' on the sets that have one. By the Hahn--Banach theorem (Part~II) they
extend to all subsets as invariant means, but no such extension can be written down explicitly.
\end{solution}

\begin{solution}{pb:bollobas}
\idea Order the ground set at random. The event ``all of $A_i$ before all of $B_i$'' has probability $1/\binom{a_i+b_i}{a_i}$, and two such
events cannot happen together, because of an element of $A_i\cap B_j$ and one of $A_j\cap B_i$.

\writeup (a) Let $X$ be the union of all the sets, and $\sigma$ a uniformly random ordering of $X$. Let $E_i$ be the event that every element
of $A_i$ comes before every element of $B_i$. Only the relative order of the $a_i+b_i$ elements of $A_i\cup B_i$ matters, and all orders
are equally likely; exactly $a_i!\,b_i!$ of them put $A_i$ first. So $P(E_i)=\frac{a_i!b_i!}{(a_i+b_i)!}$. Suppose $E_i$ and $E_j$ both occur, $i\ne j$.
Take $x\in A_i\cap B_j$ and $y\in A_j\cap B_i$. By $E_i$, $x$ (in $A_i$) comes before $y$ (in $B_i$); by $E_j$, $y$ (in $A_j$) comes before $x$ (in
$B_j$). This is absurd, so the events are disjoint, and $\sum_iP(E_i)\le1$.

(b) Let $G$ be $\tau$-critical with $\tau(G)=t+1$ and edges $e_1,\dots,e_m$. For each $i$, $G-e_i$ has a vertex cover $B_i$ of size $t$. It
does not cover $e_i$ (otherwise it would cover $G$), so $B_i\cap e_i=\varnothing$, and it covers every $e_j$, $j\ne i$. With $A_i=e_i$, (a) gives
$m\binom{t+2}2^{-1}\le1$. In $K_{t+2}$, a vertex cover misses at most one vertex, so $\tau=t+1$, and deleting the edge $uv$ allows the cover
$V\setminus\{u,v\}$ of size $t$.

\tools Volume~3, Chapter~1 (uniform random permutations); Volume~1 (the LYM inequality, which is the case $B_i=X\setminus A_i$).
\leads The inequality has a ``skew'' version, where one only assumes $A_i\cap B_j\ne\varnothing$ for $i<j$; Lovász proved it with exterior
algebra (Volume~2), a striking use of linear algebra in combinatorics. Bollobás's inequality is a basic tool in the theory of
saturated graphs and hypergraphs.
\end{solution}

\begin{solution}{pb:granularity}
\idea Boolean functions are integer-valued, so in $\{0,1\}$ coordinates their polynomial has integer coefficients; changing back to
$\pm1$ coordinates costs a factor $2$ per variable of a monomial, at most $2^d$. A coarse spectrum cannot be spread thin, and the
same integrality shows that every relevant variable carries a definite amount of influence.

\writeup (a) Let $g=\frac{1-f}2\in\{0,1\}$ and $y_i=\frac{1-x_i}2\in\{0,1\}$. As a function of $y\in\{0,1\}^n$, $g$ has a unique multilinear expansion
$g=\sum_Ta_Ty^T$, and evaluating at $y=1_V$ gives $g(1_V)=\sum_{T\subseteq V}a_T$; by Möbius inversion on subsets, $a_T=\sum_{U\subseteq T}(-1)^{|T\setminus U|}g(1_U)$,
an integer. The substitutions $x_i=1-2y_i$ and $y_i=\frac{1-x_i}2$ are affine in each variable and preserve multilinearity, so they do not
raise the degree in either direction; as expansions are unique, $\deg_yg=\deg_xg=d$, and $a_T=0$ for $|T|>d$. Now
$y^T=2^{-|T|}\prod_{i\in T}(1-x_i)$ has $x$-coefficients in $2^{-|T|}\Z\subseteq2^{-d}\Z$, so every $\hat g(S)$ is in $2^{-d}\Z$, and $\hat f(S)=1_{S=\emptyset}-2\hat g(S)\in2^{1-d}\Z$.

(b) Each non-zero $\hat f(S)^2$ is at least $4^{1-d}$, and they add up to $1$.

(c) Induction on $n$. For $n=0$, $p$ is a non-zero constant. Write $p=x_nq+r$, with $q$, $r$ multilinear in $x'=(x_1,\dots,x_{n-1})$ and
$\deg q\le k-1$. If $q=0$, then $p=r\ne0$ and the induction hypothesis applies. Otherwise, for each $x'$ with $q(x')\ne0$, at least one
of $p(x',1)=r+q$ and $p(x',-1)=r-q$ is non-zero; so $P(p\ne0)\ge\frac12P(q\ne0)\ge\frac12\cdot2^{-(k-1)}=2^{-k}$.

(d) $D_if=\sum_{S\ni i}\hat f(S)x^{S\setminus\{i\}}$ (Problem~\ref{pb:influence}) has degree at most $d-1$, is non-zero iff $i$ is relevant, and
$\operatorname{Inf}_i(f)=P(D_if\ne0)$. By (c), $\operatorname{Inf}_i(f)\ge2^{1-d}$ for relevant $i$. Since $I(f)=\sum_S|S|\hat f(S)^2\le d$, the number $R$ of relevant variables
satisfies $R\,2^{1-d}\le d$, that is $R\le d\,2^{d-1}$.

\tools{Volume~3, Chapter~1 (inclusion--exclusion, Möbius inversion on subsets); Volume~0, Chapter~3 (polynomials).}
\leads The bound $d\,2^{d-1}$ is due to Nisan and Szegedy (1994), in their study of the degree of Boolean functions as real
polynomials. Chiarelli, Hatami and Saks (2020) improved it to $O(2^d)$, which is the right order: the ``address'' function, where $d-1$
bits point to one of $2^{d-1}$ other bits and the output is the bit pointed to, has degree $d$ and $d-1+2^{d-1}$ relevant variables.
\end{solution}

\section*{Sources and further reading}

\begin{itemize}
\item A.~N.~Kolmogorov, \emph{Grundbegriffe der Wahrscheinlichkeitsrechnung}, Springer, 1933 (English:
\emph{Foundations of the Theory of Probability}, Chelsea, 1950). The axioms, in sixty pages.
\item W.~Feller, \emph{An Introduction to Probability Theory and Its Applications}, Vol.~1, 3rd ed., Wiley, 1968,
Chapters~I--V and~III.1 for the ballot theorem and the reflection principle.
\item J.~Bertrand, \emph{Calcul des probabilités}, Gauthier-Villars, 1889. The chord paradox (pp.~4--5) and the
ballot problem.
\item G.-L.~Leclerc de Buffon, \emph{Essai d'arithmétique morale}, Supplément à l'Histoire naturelle, vol.~4, 1777.
The needle; J.-É.~Barbier, \emph{Note sur le problème de l'aiguille et le jeu du joint couvert}, Journal de
Mathématiques Pures et Appliquées 5 (1860), 273--286, for the noodle.
\item E.~T.~Jaynes, \emph{The well-posed problem}, Foundations of Physics 3 (1973), 477--492. Invariance resolves
Bertrand's paradox.
\item S.~Selvin, \emph{A problem in probability} (letter), The American Statistician 29 (1975), 67; and
J.~P.~Morgan, N.~R.~Chaganty, R.~C.~Dahiya and M.~J.~Doviak, \emph{Let's make a deal: the player's dilemma},
The American Statistician 45 (1991), 284--287. Monty Hall and the role of the host's rule.
\item L.~Badger, \emph{Lazzarini's lucky approximation of $\pi$}, Mathematics Magazine 67 (1994), 83--91.
\item M.~Mitzenmacher and E.~Upfal, \emph{Probability and Computing}, 2nd ed., Cambridge University Press, 2017,
Chapter~5. Birthdays, balls into bins and hashing.
\item D.~V.~Lindley, \emph{Dynamic programming and decision theory}, Applied Statistics 10 (1961), 39--51; and
T.~S.~Ferguson, \emph{Who solved the secretary problem?}, Statistical Science 4 (1989), 282--296.
\end{itemize}
